\documentclass[12pt]{article}
\usepackage[utf8]{inputenc}
\usepackage[turkish]{babel}
\usepackage[a4paper, margin=2.5cm]{geometry}
\usepackage{hyperref}
\usepackage{booktabs}
\usepackage{tabularx}
\usepackage{amsmath}
\usepackage{amssymb}
\usepackage{titlesec}
\usepackage{xcolor}
\usepackage{listings}

% Kod blokları için renk ve biçim tanımlamaları (Premium stil)
\definecolor{codegreen}{rgb}{0,0.5,0}
\definecolor{codegray}{rgb}{0.4,0.4,0.4}
\definecolor{codepurple}{rgb}{0.58,0,0.82}
\definecolor{backcolour}{rgb}{0.96,0.96,0.96}

\lstdefinestyle{javastyle}{
    backgroundcolor=\color{backcolour},   
    commentstyle=\color{codegreen}\itshape,
    keywordstyle=\color{blue}\bfseries,
    numberstyle=\tiny\color{codegray},
    stringstyle=\color{codepurple},
    basicstyle=\ttfamily\footnotesize,
    breakatwhitespace=false,         
    breaklines=true,                 
    captionpos=b,                    
    keepspaces=true,                 
    numbers=left,                    
    numbersep=5pt,                  
    showspaces=false,                
    showstringspaces=false,
    showtabs=false,                  
    tabsize=2,
    language=Java
}
\lstset{style=javastyle}

\hypersetup{
    colorlinks=true,
    linkcolor=blue,
    filecolor=magenta,      
    urlcolor=cyan,
    pdftitle={slimeSlayer - Oyun Tasarim ve Uygulama Raporu},
}

\title{
    \vspace{-2cm}
    \textbf{\Large Okul Projesi Teknik Raporu}\\
    \vspace{0.5cm}
    \textbf{\huge ``slimeSlayer''}\\
    \vspace{0.5cm}
    \large 2D Top-Down Survival Roguelite Oyunu Tasarım ve Uygulama Detayları
}

\author{
    \textbf{Geliştirici Ekibi:}\\
    Gizem Öksüz \\
    Ali Andaç Erdaş \\
    İsmail Emre Türkbay
}
\date{\today}

\begin{document}

\maketitle
\newpage
\tableofcontents
\newpage

\section{Giriş ve Proje Hakkında}
Bu rapor; okul projesi kapsamında Java platformunun grafik arayüz kütüphaneleri (Swing ve AWT) kullanılarak harici bir oyun motoru desteği olmaksızın geliştirilen \textbf{``slimeSlayer''} oyununun mimarisini, mekaniklerini, matematiksel altyapısını ve optimizasyon süreçlerini incelemektedir.

\subsection{Oyunun Amacı}
\textbf{slimeSlayer}, oyuncunun bataklık temalı bir harita üzerinde hayatta kalma mücadelesini ele alan 2D roguelite türünde bir oyundur. Oyuncu WASD tuşları yardımıyla hareket ederken, otomatik saldıran üç temel yetenek (Ateş Topu, Döner Bıçak ve Kalkan) ile sürekli çoğalan Slime ve Golem ordularına karşı direnir. Düşmanlardan düşen deneyim kristallerini topladıkça seviye atlar, kart seçim ekranından yeteneklerini geliştirir, haritada doğan kırmızı güç meyvelerini yiyerek geçici hız kazanır. Ayrıca oyunda hasar anında uçuşan hasar sayıları, sentezlenmiş retro ses efektleri, ses seviyesi ve zorluk ayarlarının yapılabildiği interaktif bir giriş ekranı bulunur.

---

\section{Proje Ekibi ve Rol Dağılımı}
Projede yer alan ekip üyeleri ve geliştirme süreçlerindeki katkıları şu şekildedir:
\begin{itemize}
    \item \textbf{Gizem Öksüz:} Seviye atlama kart arayüzü (\texttt{SeviyeArayuzu}), slot temelli envanter arayüzü, can/xp barları tasarımı ve entegrasyonu.
    \item \textbf{Ali Andaç Erdaş:} Zemin çizim optimizasyonu, mıknatıslı kristaller, kırmızı hız meyvesi, döner bıçakların yörüngesi, hasar anında uçuşan renkli hasar sayıları (\texttt{HasarSayisi}), retro ses sentezleyici motoru (\texttt{SesSentezleyici}) ve ana giriş arayüzü (\texttt{MenuArayuzu}).
    \item \textbf{İsmail Emre Türkbay:} Düşman yapay zekası (\texttt{Dusman}), zıplama/yürüme animasyonlarının spritesheet üzerinden kesilmesi, Golem ve Hızlı düşman varyasyonları, dalga uretim algoritması ve itme/fizik motoru.
\end{itemize}

---

\section{Oyun Mekanikleri ve Kod Tasarımı}

\subsection{Karakter Hareketi ve Çapraz Hız Normalizasyonu}
8 yönlü hareket sistemlerinde, yatay ve dikey tuşlara aynı anda basıldığında ($\sqrt{2} \approx 1.414$ kat) ortaya çıkan hız artışı, \texttt{Oyuncu.java} içerisinde normalize edilmiştir:

\begin{lstlisting}[caption={Çapraz Harekette Hız Normalizasyonu}, label={list:normalizasyon}]
// Capraz harekette hizin 1.41 katina cikmasini engellemek icin normalizasyon degeri (sin(45) = 0.707)
if (hareketX != 0 && hareketY != 0) {
    hareketX *= 0.707;
    hareketY *= 0.707;
}
double toplamHiz = hiz + ekstraHiz;
x += hareketX * toplamHiz;
y += hareketY * toplamHiz;
\end{lstlisting}

\subsection{Düşman Spawn Algoritması ve Zorluk Ölçeklemesi}
\texttt{DusmanUretici.java} sınıfı, oyuncu etrafındaki ekran sınırlarının hemen dışındaki 4 kenardan (+60px ofset) düşman spawn eder. Zorluk derecesi, ana menüden seçilen \texttt{zorlukModu} çarpanı ile çarpılarak düşmanların can ve hasarlarına etki eder:
\begin{equation}
    \text{Zorluk Carpani} = \left( 1.0 + \frac{\text{Gecen Sure (ms)}}{120000.0} \right) \times \text{zorlukModu}
\end{equation}

\subsection{Mıknatıs Çekimli Tecrübe Kristalleri}
Düşmanlar elendiğinde yere tecrübe kristalleri düşer. Oyuncu 120 piksel toplama menziline girdiğinde mıknatıs çekim alanı devreye girer:
\begin{equation}
    \vec{V}_{cekme} = \vec{u}_{yon} \times \min(12.0, V_{eski} + 0.25)
\end{equation}

---

\section{Grafik, Görsel Efektler ve Çizim Yönetimi}

\subsection{Programatik Renklendirme}
Düşman varyasyonları için diskten tek bir spritesheet okunur ve \texttt{GorselYukleyici.java} içinde programatik gri tonlama algoritmasıyla gri golem bossuna dönüştürülür:

\begin{lstlisting}[caption={Programatik Gri Tonlama Filtresi}, label={list:gri}]
public static BufferedImage gorseliGriyap(BufferedImage kaynak) {
    BufferedImage sonuc = new BufferedImage(kaynak.getWidth(), kaynak.getHeight(), BufferedImage.TYPE_INT_ARGB);
    for (int x = 0; x < kaynak.getWidth(); x++) {
        for (int y = 0; y < kaynak.getHeight(); y++) {
            int rgb = kaynak.getRGB(x, y);
            int a = (rgb >> 24) & 0xff;
            int r = (rgb >> 16) & 0xff;
            int g = (rgb >> 8) & 0xff;
            int b = rgb & 0xff;
            int gri = (int) (r * 0.299 + g * 0.587 + b * 0.114);
            int yeniRgb = (a << 24) | (gri << 16) | (gri << 8) | gri;
            sonuc.setRGB(x, y, yeniRgb);
        }
    }
    return sonuc;
}
\end{lstlisting}

\subsection{Performans Optimizasyonu (Frustum Culling)}
Büyük haritalarda (3000x3000px) tüm zemini çizmek ağır performans kaybına yol açar. Projede sadece kameranın gördüğü alandaki karolar taranıp çizilir. Ayrıca ekran sınırları dışındaki düşmanlar \texttt{getClipBounds} ile elenir (\textbf{Frustum Culling}).

\subsection{Kamera Sarsıntısı (Screen Shake) ve Hasar Ekranı Flaş Efekti}
Oyuncu hasar aldığında vuruş hissini artırmak adına kamera rastgele sarsılır ve ekrana kırmızı saydam bir perde çekilir. Sarsıntı gücü ve flaşın opasitesi zamanla azalır:

\begin{lstlisting}[caption={Kamera Sarsıntısı ve Flaş Uygulaması}, label={list:sarsinti}]
// OyunPaneli.java paintComponent metodu
double nihaiKameraX = kameraX;
double nihaiKameraY = kameraY;

if (sarsintiSuresi > 0) {
    java.util.Random rand = new java.util.Random();
    nihaiKameraX += (rand.nextDouble() - 0.5) * 2.0 * sarsintiGucu;
    nihaiKameraY += (rand.nextDouble() - 0.5) * 2.0 * sarsintiGucu;
}
g2.translate(-nihaiKameraX, -nihaiKameraY);
// ... Cizimler ...
g2.translate(nihaiKameraX, nihaiKameraY); // HUD oncesi geri al

// Hasar Flas perdesi
if (hasarFlasiSuresi > 0) {
    int alfa = (int) (110.0 * ((double) hasarFlasiSuresi / 8.0));
    g2.setColor(new Color(255, 0, 0, alfa));
    g2.fillRect(0, 0, EKRAN_GENISLIGI, EKRAN_YUKSEKLIGI);
}
\end{lstlisting}

---

\section{Retro Ses Efektı Sentezleyici Motoru}

Ses sentezleyici motoru, diskteki `.wav` dosyalarının kaybolması veya işletim sisteminin ses sürücelerinin hata vermesi risklerini bertaraf etmek amacıyla tamamen **Java Sound API** (\texttt{javax.sound.sampled}) kullanılarak sıfırdan yazılmıştır. Sentezlenen 16-bit PCM dalgalar ayrı bir thread (iş parçacığı) üzerinden asenkron olarak çalınarak FPS düşüşleri engellenir.

\begin{lstlisting}[caption={Asenkron 16-bit PCM Mono Ses Çalma Modülü}, label={list:sessentezleyici}]
private static void sesCal(float[] data) {
    new Thread(() -> {
        try {
            byte[] byteData = new byte[data.length * 2];
            for (int i = 0; i < data.length; i++) {
                short val = (short) (data[i] * 32767); // Genligi 16-bit short'a donustur
                byteData[i * 2] = (byte) (val & 0xff);       // Little-endian alt byte
                byteData[i * 2 + 1] = (byte) ((val >> 8) & 0xff); // High byte
            }
            AudioFormat format = new AudioFormat(44100, 16, 1, true, false);
            SourceDataLine line = AudioSystem.getSourceDataLine(format);
            line.open(format);
            line.start();
            line.write(byteData, 0, byteData.length);
            line.drain();
            line.close();
        } catch (Exception e) {
            System.out.println("Ses calinirken hata: " + e.getMessage());
        }
    }).start();
}
\end{lstlisting}

---

\section{Önemli Kod Blokları ve Teknik Analiz}

\subsection{Vektörel Çarpışma Çözücü ve Geri Sekme (Knockback)}
Daire tabanlı çarpışmalarda iki nesnenin iç içe girmesini önlemek amacıyla çarpışan alanların yarıçap toplamı kadar nesneler zıt yönde itilir. Golem Boss'unun direnç katsayısı (\texttt{geriItmeCarpani = 0.15}) ağır kütleyi başarıyla simüle eder:

\begin{lstlisting}[caption={Vektörel Çarpışma Çözücü}, label={list:carpismafizik}]
// CarpismaDenetleyici.java
double dx = oyuncu.x - dusman.x;
double dy = oyuncu.y - dusman.y;
double mesafeKaresi = dx * dx + dy * dy;
double yariCapToplam = oyuncu.yariCap + dusman.yariCap;

if (mesafeKaresi < yariCapToplam * yariCapToplam) {
    double mesafe = Math.sqrt(mesafeKaresi);
    if (mesafe > 0) {
        double overlap = yariCapToplam - mesafe;
        double itmeX = (dx / mesafe) * overlap;
        double itmeY = (dy / mesafe) * overlap;
        
        oyuncu.x += itmeX * 0.6;
        oyuncu.y += itmeY * 0.6;
        dusman.x -= itmeX * 0.4 * dusman.geriItmeCarpani;
        dusman.y -= itmeY * 0.4 * dusman.geriItmeCarpani;
    }
}
\end{lstlisting}

\subsection{Döner Bıçakların Yörüngesi ve Hasar Cooldown Sistemi}
Döner bıçak mermileri oyuncu etrafında dairesel yörüngede döner. Denge sağlamak amacıyla düşman bazlı hasar zaman damgası takibi yapılmıştır:
\begin{equation}
    x_{bicak} = x_{oyuncu} + \cos(\theta) \times r_{yorunge}
\end{equation}
\begin{equation}
    y_{bicak} = y_{oyuncu} + \sin(\theta) \times r_{yorunge}
\end{equation}

\begin{lstlisting}[caption={Döner Bıçak Hasar Zaman Damgası Takibi}, label={list:hashmapcooldown}]
if (mermi instanceof DonerBicakMermi) {
    DonerBicakMermi dbMermi = (DonerBicakMermi) mermi;
    long suAn = System.currentTimeMillis();
    
    // Dusmana tekrar vurup vuramayacagini denetleyen metot (500 ms bekleme suresi)
    if (dbMermi.dusmanaVurabilirMi(dusman, suAn)) {
        dusman.can -= mermi.hasar;
        dbMermi.vurulanDusmaniEkle(dusman, suAn); // HashMap kaydi
        
        double mesafe = Math.sqrt(mesafeKaresi);
        if (mesafe > 0) {
            dusman.geriIt(dx / mesafe, dy / mesafe, 15.0);
        }
    }
}
\end{lstlisting}

---

\section{Kullanılan Java Yapıları ve Tercih Nedenleri}

\subsection{Nesne Yönelimli Programlama (OOP)}
\begin{itemize}
    \item \textbf{Kalıtım (Inheritance):} \texttt{GolemDusman} ve \texttt{HizliDusman} sınıfları \texttt{Dusman} sınıfından türetilmiştir.
    \item \textbf{Soyut Sınıflar (Abstract Classes):} \texttt{Silah} sınıfı soyut yapılarak silahlar için kurumsal standartlar konulmuş, \texttt{saldir()} metodu her silahta özelleştirilmiştir.
    \item \textbf{Çok Biçimlilik (Polymorphism):} \texttt{Silah} türündeki listeler içinde barındırılan nesneler döngülerde aynı imza ile çağrılmış fakat kendi metotlarını tetiklemiştir.
\end{itemize}

\subsection{Koleksiyon Yapısı (Collections Framework)}
\begin{itemize}
    \item \textbf{ArrayList:} Haritadaki dinamik varlık listelerini (mermiler, kristaller, düşmanlar) yönetmek için \texttt{ArrayList} yapısı seçilmiştir.
    \item \textbf{HashMap:} Dönen bıçak hasar cooldown takibinde düşman referanslarını anahtar (\texttt{key}), son vuruş zaman damgalarını değer (\texttt{value}) olarak saklayan \texttt{HashMap<Dusman, Long>} yapısı kullanılmıştır.
\end{itemize}

---

\section{Kaynakça}
\begin{enumerate}
    \item \textbf{Oyuncu Spritesheet Tasarımı:} \\
    \texttt{assets/Heroes99\_free 23.13.22/character\_spritesheet.png}
    \item \textbf{Canavar/Slime Animasyonları:} \\
    \texttt{assets/FreeCharactersAnimationsAssetPack 23.13.22/Monster\_Slime\_Walk-Sheet.png}
    \item \textbf{Zemin Tileset Çizimleri:} \\
    \texttt{assets/Dark\_Swamp\_Starter\_Pack\_v1.0 23.13.22/RawAssets/GroundTileset.png}
    \item \textbf{HUD Arayüz Şablonları:} \\
    \texttt{assets/Pixel UI pack 3 23.13.22/}
    \item \textbf{Referans Dokümantasyonlar:} \\
    Oracle Java Sound API - Java Sampled Sound References.
\end{enumerate}

\end{document}
